\section{Appendix to Section \ref{sec:private}}\label{app:privateproofs}

\subsection{Effort choice}\label{app:effort}

For a given project type \(j\), the entrepreneur solves
\begin{equation}
    \max_{x \in [0,1]}
    \left\{
        xR_j(q)-\frac{x^2}{2}
    \right\}.
\end{equation}
The derivative is
\begin{equation}
    \frac{\partial}{\partial x}
    \left(
        xR_j(q)-\frac{x^2}{2}
    \right)
    =
    R_j(q)-x.
\end{equation}
Hence the unconstrained optimum is \(x=R_j(q)\).
Imposing the constraint \(x \in [0,1]\) yields
\begin{equation}\label{eq:xgeneral}
    x_j^*(q)=\min\{R_j(q),1\},
\end{equation}
since \(R_j(q)>0\).
Under Assumption \ref{ass:interior}, \(R_j(q)\in(0,1)\), so \eqref{eq:xgeneral} reduces to \eqref{eq:xjstar}.

Substituting \(x_j^*(q)=R_j(q)\) into the entrepreneur's objective yields
\begin{equation}
    \max_{x \in [0,1]}
    \left\{
        xR_j(q)-\frac{x^2}{2}
    \right\}
    =
    \frac{R_j(q)^2}{2},
\end{equation}
which gives \eqref{eq:Uj}.

\subsection{Project choice without the single-crossing restriction}\label{app:projchoice}

Define the private payoff difference
\begin{equation}
    D(q)\equiv R_B(q)-R_I(q).
\end{equation}
Using the definition of \(R_j(q)\), this difference is
\begin{equation}\label{eq:Dqpiece}
    D(q)
    =
    \begin{cases}
        \pi_E^D(B)-\pi_E^D(I), &
        q < \bar q_B, \\
        \pi_E^D(B)-\pi_E^D(I)+\beta(q\Delta_B-k), &
        \bar q_B \le q < \bar q_I, \\
        \pi_E^D(B)-\pi_E^D(I)+\beta q(\Delta_B-\Delta_I), &
        q \ge \bar q_I.
    \end{cases}
\end{equation}
Because \(\pi_E^D(I)>\pi_E^D(B)\), we have \(D(q)<0\) for all \(q<\bar q_B\).
In the intermediate region \([\bar q_B,\bar q_I)\), the buyout-oriented project is preferred if and only if
\begin{equation}\label{eq:qtilde}
    q \ge \tilde q
    \equiv
    \frac{\pi_E^D(I)-\pi_E^D(B)+\beta k}{\beta \Delta_B}.
\end{equation}
Once both projects are merger-feasible, the buyout-oriented project is preferred if and only if
\begin{equation}\label{eq:qhatappendix}
    q \ge \hat q
    \equiv
    \frac{\pi_E^D(I)-\pi_E^D(B)}{\beta(\Delta_B-\Delta_I)}.
\end{equation}

Assumption \ref{ass:singlecross} implies \(\hat q \ge \bar q_I\).
This is equivalent to
\begin{equation}
    \pi_E^D(I)-\pi_E^D(B)
    \ge
    \beta k\left(\frac{\Delta_B}{\Delta_I}-1\right),
\end{equation}
which in turn implies \(\tilde q \ge \bar q_I\).
Hence no switching occurs in the intermediate region \([\bar q_B,\bar q_I)\), and the equilibrium project choice collapses to the single-threshold characterization in Proposition \ref{prop:projectchoice}.

\subsection{Proof of Proposition \ref{prop:projectchoice}}

By \eqref{eq:Dqpiece}, the payoff difference \(D(q)=R_B(q)-R_I(q)\) is weakly increasing in \(q\), and strictly increasing on any interval where at least one project is merger-feasible.

For \(q<\bar q_I\), Assumption \ref{ass:singlecross} and Appendix \ref{app:projchoice} imply \(D(q)<0\), so project \(I\) is privately optimal.

For \(q\ge \bar q_I\), both projects are merger-feasible and
\begin{equation}
    D(q)
    =
    \pi_E^D(B)-\pi_E^D(I)+\beta q(\Delta_B-\Delta_I).
\end{equation}
Because \(\Delta_B>\Delta_I\), this expression is strictly increasing in \(q\), and it equals zero at
\begin{equation}
    q=\hat q
    =
    \frac{\pi_E^D(I)-\pi_E^D(B)}
         {\beta(\Delta_B-\Delta_I)}.
\end{equation}
Assumption \ref{ass:singlecross} ensures \(\hat q \in [\bar q_I,1)\).
Therefore \(D(q)<0\) for \(q<\hat q\), \(D(q)=0\) at \(q=\hat q\), and \(D(q)>0\) for \(q>\hat q\), proving Proposition \ref{prop:projectchoice}. \qed

\subsection{Proof of Proposition \ref{prop:entrymonotonicity}}

By \eqref{eq:Rbarpiece},
\begin{equation}
    \bar R(q)
    =
    \begin{cases}
        \pi_E^D(I), &
        q < \bar q_I, \\
        \pi_E^D(I)+\beta(q\Delta_I-k), &
        \bar q_I \le q < \hat q, \\
        \pi_E^D(B)+\beta(q\Delta_B-k), &
        q \ge \hat q.
    \end{cases}
\end{equation}
Differentiating on each interval yields
\begin{equation}
    \frac{d\bar R(q)}{dq}
    =
    \begin{cases}
        0, & q < \bar q_I, \\
        \beta \Delta_I, & \bar q_I < q < \hat q, \\
        \beta \Delta_B, & q > \hat q,
    \end{cases}
\end{equation}
which proves part (i).

Part (ii) follows from \(x^*(q)=\bar R(q)\).

For part (iii), recall that under Assumption \ref{ass:interior},
\begin{equation}
    n(s,q)=s+\frac{\bar R(q)^2}{2}.
\end{equation}
Differentiating with respect to \(q\) gives
\begin{equation}
    \frac{\partial n(s,q)}{\partial q}
    =
    \bar R(q)\frac{d\bar R(q)}{dq}.
\end{equation}
Because \(\bar R(q)>0\) and \(d\bar R(q)/dq \ge 0\), entrant mass is weakly increasing in \(q\). \qed

\subsection{Proof of Corollary \ref{cor:subsidy}}

From \eqref{eq:Uj},
\begin{equation}
    U_j(f;s,q)=s-f+\frac{R_j(q)^2}{2}.
\end{equation}
Since the term \(s\) enters identically for both projects, it does not affect the ranking of \(U_I\) and \(U_B\), and therefore it does not affect private project choice.
Because effort depends only on the chosen project's continuation value \(R_{j^*(q)}(q)\), effort is likewise unaffected.

From \eqref{eq:entrycutoff},
\begin{equation}
    f^*(s,q)=s+\frac{\bar R(q)^2}{2},
\end{equation}
so
\begin{equation}
    \frac{\partial f^*(s,q)}{\partial s}=1.
\end{equation}
Under the uniform distribution,
\begin{equation}
    n(s,q)=f^*(s,q),
\end{equation}
hence
\begin{equation}
    \frac{\partial n(s,q)}{\partial s}=1.
\end{equation}
This proves Corollary \ref{cor:subsidy}. \qed

\subsection{A note on type-specific approval or exit-realizability shifts}\label{app:typespecificq}

The main text uses a single reduced-form parameter \(q\) in order to keep the private equilibrium one-dimensional.
The underlying mechanism does not require that simplification.
Suppose instead that project \(j\) faces its own approval or exit-realizability parameter \(r_j\in[0,1]\).
Continuation values become
\begin{equation}
    R_j(r_j)=\pi_E^D(j)+\beta[r_j\Delta_j-k]_+,
    \qquad j\in\{I,B\}.
\end{equation}
The entrepreneur then chooses the buyout-oriented project if and only if
\begin{equation}
    R_B(r_B)\ge R_I(r_I).
\end{equation}
Thus the composition mechanism survives whenever the effective exit shift for project \(B\) is sufficiently strong relative to its lower stand-alone payoff.
The baseline restriction \(r_I=r_B=q\) is adopted only to isolate that mechanism with a single policy or exit-environment shifter.

\subsection{General distributions and the scope of Appendix B}\label{app:generalG}

This appendix shows that the main private-equilibrium results do not rely on the uniform distribution beyond the closed-form expression for entrant mass.

Let the fixed entry cost \(f\) be drawn from a continuously differentiable distribution \(G\) with density \(g\), supported on an interval containing all relevant cutoffs.
Define
\begin{equation}
    \Phi(q)\equiv
    \max_{j \in \{I,B\}}
    \left\{
        \max_{x \in [0,1]}
        \left[
            xR_j(q)-\frac{x^2}{2}
        \right]
    \right\}.
\end{equation}
Under Assumption \ref{ass:interior},
\begin{equation}
    \Phi(q)=\frac{\bar R(q)^2}{2}.
\end{equation}
The entrepreneur enters if and only if
\begin{equation}
    f \le s+\Phi(q),
\end{equation}
so the mass of entrants is
\begin{equation}
    n_G(s,q)=G\bigl(s+\Phi(q)\bigr).
\end{equation}
Differentiating gives
\begin{align}
    \frac{\partial n_G(s,q)}{\partial s}
    &= g\bigl(s+\Phi(q)\bigr), \\
    \frac{\partial n_G(s,q)}{\partial q}
    &= g\bigl(s+\Phi(q)\bigr)\Phi'(q).
\end{align}
Hence the qualitative conclusions of Proposition \ref{prop:entrymonotonicity} remain unchanged:
whenever \(\Phi'(q)\ge 0\), a more permissive merger policy weakly increases entry.

Appendix B studies the private equilibrium under a general entry-cost distribution \(G(f)\).
Because the revised welfare analysis in Sections \ref{sec:subsidy} and \ref{sec:decentralized}
explicitly integrates realized fixed costs and the fiscal subsidy cost, welfare under general \(G\)
is handled in the planner-specific appendices rather than through a separate reduced-form shortcut
here. In particular, the relevant planner-\(p\) welfare object takes the form
\begin{equation}
    W_j^p(s,q)
    =
    \int_0^{c_j(s,q)}
    \left[
        R_j(q)B_j^p(q)-\frac{R_j(q)^2}{2}-f-\tau s
    \right]\,dG(f),
\end{equation}
with
\begin{equation}
    c_j(s,q)=s+\frac{R_j(q)^2}{2}.
\end{equation}
Accordingly, Appendix B contains no separate reduced-form welfare shortcut beyond the
private-equilibrium objects derived above.
